% Prologue of Verifiable Intelligence.
\specialchapter{Prologue}{The weak verifier}{V,P,V,P,V,P,D}\label{ch:prologue}
\begin{refsection}[chapters/00-prologue.bib]
\vifloatprefix{P}

In 1985 László Babai named the two players of a new kind of proof after a king and a wizard. Arthur, the king, can toss coins and do a little arithmetic, and he must decide whether a claim is true. Merlin, the wizard, knows everything, but he cannot be trusted: he wants Arthur to accept his claim, whether or not it is true. Babai asked which claims Merlin can prove to Arthur, and in his paper the claims were about groups of matrices too large for Arthur to list \parencite{Babai1985}. The same year Shafi Goldwasser, Silvio Micali and Charles Rackoff gave a different definition of the same idea, with a verifier who keeps its coins secret and shows the prover only messages computed from them, and asked a second question beside the first: how much does the verifier learn from a proof beyond the fact that the claim is true \parencite{GoldwasserMicaliRackoff1989}? The journal versions of the two papers, Babai's written with Shlomo Moran \parencite{BabaiMoran1988}, shared the first Gödel Prize in 1993.\sidenote{The Gödel Prize is awarded for outstanding papers in theoretical computer science. The prize for the PCP theorem followed in 2001, and Irit Dinur's combinatorial proof of it received the prize in 2019.}

Merlin was a figure of speech. The prover of an interactive proof is all-powerful because the theory wants to know what the verifier can check \emph{whatever} the prover does; it was an idealization, not a description of a machine. The question the theory asked was nonetheless a practical one: \emph{what can a weak verifier learn from a strong prover that may lie?} Over the following decade the answers arrived in a rush, several of them by electronic mail within weeks of one another \parencite{Babai1990}. A verifier that may ask questions and flip coins can check every claim that a computer could decide in polynomial space, including claims about games far too large to search. A verifier that questions two provers in separate rooms can check claims about computations of exponential length. A verifier that reads a few bits of a proof, chosen at random, can check every claim that has a short proof at all, provided the proof is written in a suitable code. A prover can convince a verifier that a claim is true while revealing nothing else about why.

\addsec{The oversight problem}

Thirty-three years after Arthur and Merlin, Geoffrey Irving, Paul Christiano and Dario Amodei proposed a method for supervising AI systems whose answers people cannot check: train two copies of a system to debate each other in front of a human judge, rewarding whichever copy the judge believes \parencite{IrvingEtAl2018}. They pointed out that the question had a precedent. A judge who can follow a short argument but cannot find the answer is a verifier, the debaters are provers, and a judge who weighs a debate between optimal players can decide every question in $\cc{PSPACE}$, the class that interactive proofs reach, where a judge who weighs a single answer reaches only $\cc{NP}$. Two years earlier, a list of concrete problems in AI safety had named the difficulty \emph{scalable oversight}: how to make a system behave as intended when the objective is too expensive to evaluate often during training \parencite{AmodeiEtAl2016}. Later work extended the term to tasks on which the system outperforms the people who supervise it.\sidenote{The same list includes avoiding negative side effects, avoiding reward hacking, safe exploration, and robustness to distributional shift. Scalable oversight is the one this book is about.} Experimental work has since tried to measure progress on it, for example by testing whether people who are not experts in a task, working with a model that is better at the task than they are, can do better than either alone, with experts supplying the right answers \parencite{BowmanEtAl2022}.

The difficulty grows with capability. As of 2025, language models write long programs and changes to large codebases, produce mathematical proofs that a proof assistant can check \parencite{HubertEtAl2026}, propose plans and experiments, and forecast events; the people who oversee them can check some of these outputs, some of the time, and fewer as the outputs grow. A person who could produce the answer can check a model by comparing answers. A person who cannot must check it some other way, and must do so knowing that the model may be wrong in ways that pass the check: by mistake, because training rewarded answers that look right, or because the model is pursuing a goal of its own. This is the oversight problem in the form this book studies:

\begin{oversight}
A model makes a claim and may support it with whatever evidence it chooses. The overseer can afford a small amount of computation, can ask questions, can sample at random, and can consult a human judgment now and then, but cannot find the answer itself. The overseer must accept the claim when it is true, at least when the model argues for it honestly, and must reject it when it is false, whatever the model does.
\end{oversight}

\noindent Read as mathematics, this is the definition of an interactive proof. Read as engineering, it is a list of things that are hard to guarantee: that the overseer's questions cannot be predicted, that the model's copies cannot coordinate, that the honest argument is one the model can afford to make, that the human judgment consulted is the judgment wanted. The book takes both readings seriously. Each chapter teaches one protocol from the theory, proves what the theory proves about it, and then asks what the proof assumes that oversight cannot.

\addsec{A first protocol}

A small example shows the shape of everything that follows. A model has been asked to simplify an algebraic expression with thousands of terms and integer coefficients, and it returns a much shorter one that it claims is equal. Expanding both expressions to compare them would take longer than the overseer can afford. Evaluating them is cheap: after the model has answered, the overseer draws a large prime $p$ at random and a number $r$ uniformly from $\{0,1,\dots,p-1\}$, evaluates both expressions at $r$ with arithmetic modulo $p$, and accepts if the two values agree.\sidenote{With several variables, the overseer draws one random value for each. The bound $d/p$ below holds for every number of variables, by the lemma of Schwartz and Zippel (\cref{lem:toolkit:schwartzzippel}).}

If the claim is true, the two expressions are the same polynomial, and the values agree at every point. If it is false, their difference is a nonzero polynomial with integer coefficients, and the lie passes only if $p$ divides all of them or $r$ is a root of the difference modulo $p$. A nonzero integer of $B$ bits has fewer than $B/60$ prime factors above $2^{60}$, while more than $2.7\times10^{16}$ primes lie between $2^{60}$ and $2^{61}$, so a prime drawn uniformly from that range divides a given nonzero coefficient with probability below $B\cdot10^{-18}$.\sidenote{This is why the prime is drawn after the claim is made. A model that knew $p$ in advance could return an expression that differs from the true one by $p$ times another polynomial, and it would pass every evaluation modulo $p$.} If both expressions have degree at most $d$, a difference that is nonzero modulo $p$ has at most $d$ roots, so $r$ is one of them with probability at most $d/p$. With $d=10^4$ and coefficients of fewer than a million bits, a false simplification survives the check with probability below $10^{-12}$. The overseer did two evaluations, while finding the simplification, or even expanding the original, may have taken far longer.

This tiny protocol already answers the three questions that every chapter of the book will ask. Is honesty the model's best policy? Here the model's only move is the claim, and a false claim is accepted with the tiny probability just computed while a true one is always accepted, so a model that wants its claim accepted does best by making a true one. What does the overseer spend? Two evaluations, a few dozen primality tests to find $p$, and the random bits that choose $p$ and $r$. How often is a lie caught? With probability above $1-10^{-12}$ in the example, \emph{whatever} false expression the model chose, because the randomness that catches it belongs to the overseer and is drawn after the claim is made. The rest of the book is about what happens when the claim is not an algebraic identity: when the overseer must ask questions, when it can read only a few bits of the evidence, when it must set two models against each other, when the model is trained rather than adversarial, and when the evidence must stay secret.

\addsec{The three questions}

Every chapter has a section called \emph{The three questions}, which answers them for its protocol and records the answers on a card. Each question has a precise reading.

\paragraph{Is honesty an equilibrium?} A protocol is a game. In the simplest reading the prover is paid when the verifier accepts. \emph{Completeness} says that on a true claim the prescribed (``honest'') strategy is accepted with probability at least $c$, usually $1$; \emph{soundness} says that on a false claim no strategy at all is accepted with probability more than $s$. If the prover first chooses which answer to assert and the protocol can prove either answer, then asserting the truth and following the honest strategy earns at least $c$, and every lie earns at most $s$. When $c=1$, honesty is therefore a best response, and when $c>s$, every lie earns strictly less than honesty. Debate (\cref{ch:debate}) is a zero-sum game in which honesty is a minimax strategy; prover–verifier games (\cref{ch:pvg}) ask for equilibria in which the verifier moves first; rational proofs (\cref{ch:rational}) ask for strictly dominant truthful reports; and protocols with several provers (\cref{ch:crossexam}) ask about profiles of strategies and collusion. The chapters name their game and equilibrium notion, and ask what the honest strategy costs, which matters most for AI. An equilibrium that the honest prover cannot compute protects no one: a theorem that the truth would win against an optimal liar says nothing about a model that can afford only part of the truth.

\begin{figure}[t]
\begin{fullwidth}
\centering
\begin{tikzpicture}[x=1mm,y=1mm,font=\small]
  \def\prolrung#1#2#3#4#5{%
    \fill[#5] (7,#1) rectangle (170,#1+9.2);
    \node[anchor=west,text=prover,font=\normalsize] at (9,#1+4.6) {#2};
    \node[anchor=west,text width=80mm,align=left] at (60,#1+4.6) {#3};
    \node[anchor=east,text=verifier,font=\scriptsize\scshape] at (168,#1+4.6) {#4};}
  \prolrung{60}{$\cc{MIP}^*=\cc{RE}$}{two provers apart, sharing entanglement}{chapter 8}{prover!6}
  \prolrung{50}{$\cc{MIP}=\cc{NEXP}$}{two provers questioned apart, who cannot confer}{chapter 8}{prover!10}
  \prolrung{40}{$\cc{IP}=\cc{PSPACE}$}{polynomially many rounds of questions; debate with optimal play}{chapters 4, 5, 7}{prover!14}
  \prolrung{30}{$\cc{AM}$, $\cc{MA}$}{a constant number of rounds of coins and answers}{chapter 4}{prover!18}
  \prolrung{20}{$\cc{NP}=\cc{PCP}(O(\log n),O(1))$}{one message; or a few bits of it, if it is encoded}{chapters 1, 2}{prover!22}
  \prolrung{10}{$\cc{P}$, $\cc{BPP}$}{the verifier alone; or, with an efficient honest prover, a computation far longer than the verifier could run}{chapters 6, 9}{prover!26}
  \draw[verifier,line width=0.7pt,-{Stealth[length=2mm]}] (4.2,11) -- (4.2,69) node[midway,rotate=90,above,font=\scriptsize\scshape,text=verifier] {more to check with};
\end{tikzpicture}
\end{fullwidth}
\caption{What a weak verifier can check, by the resources it is allowed. Each class is the set of claims that a polynomial-time verifier can check with the resources on its rung; higher rungs contain the lower ones, except that whether $\cc{BPP}\subseteq\cc{NP}$ is open. The bottom rung also holds protocols that bound the honest prover's work and check a long computation in $\cc{P}$ in much less time than it takes: delegation, for computations of small depth (\cref{ch:delegation}), and doubly-efficient debate (\cref{ch:efficientdebate}).}\label{fig:prologue:ladder}
\end{figure}

\paragraph{What does the verifier spend?} Time, random bits, rounds of interaction, bits of communication, queries to a proof or a program, samples from a distribution, and calls to a human judge. The chapters state each as a function of the size of the input and of the error and security parameters. For AI oversight the scarcest of these is the last: human attention.

\paragraph{How often is a lie caught?} The probability $1-s$ that a false claim is rejected, per run and after $t$ independent repetitions, and against which class of cheaters. The answers differ in kind. Some protocols catch every prover, however powerful (information-theoretic soundness). Some catch only provers that cannot find collisions of a hash function or forge a signature (computational soundness). Some catch only provers that cannot communicate with each other. A rational proof does not catch a lie at all: it makes the lie cost money in expectation. And a protocol trained into neural networks is measured against the cheaters it saw during training, which is a different and weaker kind of guarantee. A chapter's answer to the third question always says which kind it is.

\addsec{The ladder}

The theory's answers can be arranged by the resources the verifier is allowed (\cref{fig:prologue:ladder}). On the bottom rung the verifier works alone. One rung up, it receives one message, a certificate, and checks it; these are the claims in $\cc{NP}$ (\cref{ch:certificates}), and the PCP theorem says that the verifier can check them by reading a constant number of bits of a suitably encoded certificate (\cref{ch:spotchecks}). With interaction and random questions the verifier can check every claim in $\cc{PSPACE}$, the claims decidable with a polynomial amount of memory, whose decision may take exponential time (\cref{ch:interaction,ch:sumcheck}); debate between two provers with optimal play reaches the same rung (\cref{ch:debate}). Two provers questioned in separate rooms reach $\cc{NEXP}$, nondeterministic exponential time (\cref{ch:crossexam}), and two provers who share quantum entanglement but still cannot communicate reach every claim whose truth a computer could confirm by an unbounded search, such as the claim that a given program halts, which no computer can decide \parencite{JiEtAl2021}.

The ladder measures power, and power is not what oversight most needs. Every protocol on the upper rungs asks the honest prover to do enormous work: to prove a $\cc{PSPACE}$ claim, the honest prover of the classical protocol must compute sums over exponentially many terms, and the honest debater must play a game perfectly. A model that could do that would not need overseeing on that claim, and a model that cannot do it cannot follow the honest strategy. The theory has a second axis for this reason, which asks for protocols in which the honest prover's work is close to the work of solving the problem, and the verifier's work is much less (\cref{ch:delegation,ch:efficientdebate}). In the language of the three questions, the ladder answers the third question for an all-powerful prover, and the second axis returns to the first.

\addsec{A dictionary}

The two vocabularies match closely, and the places where they do not match are where most of the interesting problems are.

\begin{fullwidth}
\begin{dictionary}
verifier $\V$, polynomial time & the overseer: a person with limited time, a weaker model, a test harness, an auditor \\
prover $\Prov$, unbounded & the model whose output is checked; powerful but bounded \\
cheating prover $\Prov^*$ & a model that is wrong, or trained to look right, or deceptive \\
the claim $x\in L$ & ``this patch is safe'', ``this plan is optimal'', ``this output came from that model'' \\
completeness $c$ & a model that argues honestly for a true claim is believed \\
soundness error $s$ & the chance that a false claim survives oversight \\
the verifier's private coins & secret, randomized tests the model cannot predict \\
a query to a proof oracle & a spot check of one step of a long argument \\
two provers who cannot confer & two instances of a model questioned in separate contexts \\
an oracle call to $\party{H}$ & one question put to a human \\
zero knowledge & an audit that reveals nothing but its verdict \\
an equilibrium & what training is hoped to find
\end{dictionary}
\end{fullwidth}

Three gaps run through the book. The first is in the verifier. The theory's verifier is a polynomial-time machine that never makes a mistake in arithmetic and judges every step it checks correctly; a human overseer is slow, inconsistent, persuadable, and has systematic blind spots. Protocols that consult a human only as an oracle for small judgments (\cref{ch:efficientdebate}) make this explicit, and they also make the human's errors part of what the protocol certifies as true. The second gap is in the prover. The theory's prover is either honest or a worst-case adversary with unlimited computation; a model is neither. Its computation is bounded, which can help the verifier, since a bounded prover cannot break a hash function (\cref{ch:training}), and can hurt it, since a bounded honest prover may be unable to follow the honest strategy (\cref{ch:sumcheck,ch:debate}). It is the product of training, which looks for high reward rather than for an equilibrium of a game, so the chapters on incentives ask what payment makes the truth worth telling (\cref{ch:rational}), what training is likely to find (\cref{ch:pvg}), and how a learned model can be checked with fewer labels than it took to learn (\cref{ch:pacverify}). The third gap is in the claim. The theory checks formal statements, and oversight cares about intentions. A protocol that certifies a formal statement certifies nothing about whether the statement says what the overseer meant (\cref{ch:certificates}).

\addsec{How to read this book}

The book has fifteen chapters in five parts. Part I (Fixed Provers) studies provers that are fixed objects: certificates, proofs that the verifier spot-checks, and programs that the verifier queries. Part II (Conversation) studies a single prover that answers questions, up to $\cc{IP}=\cc{PSPACE}$ and the delegation of computation to an efficient prover. Part III (Competition) studies several provers: debate, cross-examination, and debate in which the honest side must be efficient. Part IV (Incentives) studies provers that are trained or paid rather than adversarial: rational proofs, prover–verifier games, and the verification of learned models. Part V (Secrets and Limits) studies claims about objects the prover will not show, with zero knowledge and proofs of training, and ends with backdoors that no audit can detect. An epilogue collects the answers to the three questions across the book, and an appendix collects the tools from complexity theory, probability, algebra, game theory and cryptography that several chapters use.

Every chapter has the same anatomy. \emph{The oversight problem} states a situation in which a person or a weaker system must check a stronger one, and the proposals that have been made for it. \emph{The protocol} defines the protocol completely and writes it out as a transcript, with each step labeled by the party that sends a message or computes. \emph{The theory} proves the protocol's main theorem completely, so that a careful reader can check every line, and states further results with sketches or heuristic arguments, labeled as such. \emph{Running the protocol} runs it in a script against an honest prover and against cheating provers, and compares the measured acceptance rates with the theorem. \emph{The three questions} answers them and records the answers on a protocol card. \emph{What the model leaves out} says, in three paragraphs, what the theory idealized away, where its guarantees break, and what is open. Exercises, notes, and references follow; the notes end by translating the notation of the main sources into the book's.

\begin{figure}[t]
\begin{fullwidth}
\inlineband{P,V,P,V,P,V,P,V,D}

\vspace{1mm}
\inlineband{S,Q,Q,Q,D}

\vspace{1mm}
\inlineband{P,A,P,A,P,A,H,D}
\end{fullwidth}
\caption{Three bands. The upper rail is the prover's and the lower rail the verifier's; time runs to the right. Top: a conversation, as in the sumcheck protocol, with messages from the prover (indigo) answered by random challenges from the verifier (ochre), ending in a verdict. Middle: a proof written out in full, of which the verifier reads three bits. Bottom: a debate between two provers (indigo and teal), in which the judge consults a human once (the sage diamond) before deciding.}\label{fig:prologue:bands}
\end{figure}

The foot of each chapter's opening page carries a band like those of \cref{fig:prologue:bands}: the message pattern of the chapter's protocol, drawn on two rails. A heavy indigo stroke is a message from the prover, a light ochre stroke a message or a random challenge from the verifier, a teal stroke a message from a second prover or debater. A row of cells on the prover's rail is a proof written out in full, an ochre cell one that the verifier reads, and a dark cell one that no query reaches. A small box on a prover's message is a commitment, a sealed message the prover cannot change; a sage diamond below the verifier's rail is a question put to a human, and sage dots are samples drawn from the world; a ring on an ochre stroke is a payment; and the ochre disc is the verdict. The title page carries a string of the kind that the verifier of \cref{ch:spotchecks} reads: the $2048$ bits of the Hadamard encoding of the number $1985$. Three of them are ringed, at positions $x$, $y$ and $x+y$, the bits that one run of the linearity test of that chapter reads; they are $1$, $1$ and $0$, and the test accepts.

The scripts behind the figures are in \texttt{figures/scripts}. They need Python with NumPy and SciPy, run on a laptop in a minute or two, share a small kit, \texttt{protokit.py}, and print every number that the text quotes. They are meant to be run and changed: the quickest way to understand a soundness bound is to write the cheating prover that comes closest to it.

The prerequisites are undergraduate discrete mathematics, probability, linear algebra, and algorithms, at the level of polynomial time and reductions. No course in complexity theory or cryptography is assumed; each chapter defines what it uses, and the appendix collects the shared tools, with proofs of those whose proofs are short. The three chapters of Part I do not depend on one another, except that \cref{ch:consistency} uses the linearity test of \cref{ch:spotchecks} once. Part II should be read in order. Part III, best read in order, needs \cref{ch:sumcheck}, and cross-examination also \cref{ch:spotchecks}. Part IV needs the definitions of \cref{ch:interaction}, and its first chapter also the counting protocol of \cref{ch:sumcheck}. Part V needs \cref{ch:interaction}, and its second and third chapters need, respectively, the spot checks of \cref{ch:spotchecks} and the self-correction of \cref{ch:consistency}.

\notesandrefs

Babai's paper \parencite{Babai1985} and the journal version with Moran \parencite{BabaiMoran1988} introduced Arthur–Merlin games, in which the verifier's coins are public; Goldwasser, Micali and Rackoff \parencite{GoldwasserMicaliRackoff1989} introduced interactive proofs with private coins and zero knowledge. Goldwasser and Sipser showed that the two models have the same power \parencite{GoldwasserSipser1986}, as \cref{ch:interaction} explains. Babai's essay on the events of 1989 and 1990 \parencite{Babai1990} tells how $\cc{IP}=\cc{PSPACE}$ and its relatives were found, in a participant's short account of the period. The textbooks of \textcite{AroraBarak2009} and \textcite{Goldreich2008} cover interactive proofs, PCPs and their relatives; \textcite{Thaler2022} covers the interactive proofs and arguments used in practice; \textcite{Goldreich2001} covers zero knowledge.

\Textcite{IrvingEtAl2018} proposed debate and drew the analogy with interactive proofs; \textcite{AmodeiEtAl2016} named scalable oversight among the concrete problems of AI safety; \textcite{BowmanEtAl2022} proposed a way of measuring progress on it with human participants. Later work that the book discusses in detail appears in the chapters: doubly-efficient debate in \cref{ch:efficientdebate}, prover–verifier games in \cref{ch:pvg}, interactive proofs for verifying machine learning in \cref{ch:pacverify}, proofs about secret models in \cref{ch:zeroknowledge,ch:training}, and undetectable backdoors in \cref{ch:backdoors}.

The first protocol of this prologue, testing a polynomial identity at a random point, is due independently to \textcite{DeMilloLipton1978}, \textcite{Zippel1979} and \textcite{Schwartz1980}.

\printchapterbib
\end{refsection}
\vifloatrestore
